\chapter{The Score as Posterior Expectation: Belief Propagation and the
Kalman Filter}\label{ch:bp}

Chapter~\ref{ch:gaussian} computed the joint score by inverting a
$K \times K$ matrix; the score--posterior identity
(Section~\ref{sec:tweedie}) says the same object is a posterior
expectation. This chapter takes the Bayesian route seriously and rebuilds
the score by \emph{message passing}: local computations on the graph of
the posterior. It answers RQ2 and RQ3, and it is deliberately the most
self-contained chapter of the thesis: the graphical-model formalism ---
Markov random fields, factor graphs, the sum--product algorithm --- is
developed here, where it is used, and in the depth its central role
deserves. The plan: first the structure of the posterior (a tree), then
the sum--product algorithm and its exactness, with the honest accounting
of what ``linear cost'' does and does not mean for continuous variables;
then the Gaussian case, where the messages close algebraically --- no
central-limit argument anywhere --- and the sweeps turn out to be,
update by update, the classical Kalman filter and smoother; finally the
Gaussian closure on genuinely non-Gaussian chains, measured against exact
grid quadrature. The per-node closures of statistical physics (AMP, the
TAP equations), their exact breakdown time on the chain, and the critical
coupling $\alpha_c = \sqrt{2} - 1$ are developed in
Appendix~\ref{app:amp}: the results are complete and audited, but they
sit off the main line of the argument.

\section{The posterior and its graph}\label{sec:bp-problem}

Fix the setting of Chapters~\ref{ch:gaussian}--\ref{ch:laplace}: a Markov
chain prior $P_0(a) = p_0(a_0) \prod_k M(a_{k+1} | a_k)$, per-frame OU
evidence, and the posterior
\begin{equation}
  p(a \,|\, x) \;\propto\;
  \underbrace{\psi_{\mathrm{pr}}(a_0)
  \prod_{k=0}^{K-2} \psi^{(k)}_{\mathrm{tr}}(a_k, a_{k+1})}_{\text{chain prior}}
  \;
  \underbrace{\prod_{k=0}^{K-1} \psi^{(k)}_{\mathrm{ob}}(a_k)}_{\text{evidence}},
  \qquad
  \psi^{(k)}_{\mathrm{ob}}(a_k) = \Nd\big(x_k;\, e^{-t}a_k,\, \Deltat\big).
  \label{eq:bp-posterior}
\end{equation}
Structurally this is a \emph{hidden Markov model}
\citep{rabiner1989tutorial}: a latent chain observed through independent
per-frame noise, and by the score--posterior identity the $k$-th score
component needs exactly the smoothing marginal $\E[a_k | x]$ of
\eqref{eq:bp-posterior}. Before computing, we make the structure of this
object precise, in the language built for it.

\subsection{Markov random fields}\label{sec:mrf}

\begin{definition}[Markov random field]
Let $G = (V, \mathcal{E})$ be an undirected graph with one random variable
$x_i$ per node. The collection $x = (x_i)_{i \in V}$ is a \emph{Markov
random field} (MRF) if it satisfies the local Markov property: for every
node $i$,
\begin{equation}
  p\big(x_i \,\big|\, x_{V \setminus \{i\}}\big)
  \;=\; p\big(x_i \,\big|\, x_{\partial i}\big),
  \label{eq:local-markov}
\end{equation}
where $\partial i$ is the set of neighbours of $i$ in $G$.
\end{definition}

The fundamental structure theorem connects this conditional-independence
definition to the energy picture of Chapter~\ref{ch:diffusion}: by the
Hammersley--Clifford theorem \citep[proved in][]{besag1974spatial}, a
strictly positive density is an MRF on $G$ if and only if it factorises
over the \emph{cliques} of $G$,
\begin{equation}
  p(x) \;=\; \frac{1}{Z} \prod_{C \in \mathrm{cliques}(G)} \psi_C(x_C)
  \;=\; \frac{1}{Z}\, \exp\Big[ -\sum_C E_C(x_C) \Big] .
  \label{eq:hammersley-clifford}
\end{equation}
Conditional independence, graphical locality, and additive local energies
are one and the same thing. The Ising model \eqref{eq:ising} is the
textbook example; the decisive example for us is one-dimensional: the
chain factorisation \eqref{eq:chain-factorisation} makes a Markov chain an
MRF on a line, with pairwise nearest-neighbour energies. For the Gaussian
AR(1) chain each energy term is quadratic, so the coupling matrix of the
quadratic form connects only neighbours --- this is the
Hammersley--Clifford reading of the tridiagonal precision of
Chapter~\ref{ch:gaussian} (Remark~\ref{rmk:g-markov-sparsity}), now with
its proper name. And because the likelihood factors of
\eqref{eq:bp-posterior} touch one variable each, conditioning on the
observations does not enlarge any clique: \emph{the posterior is again an
MRF on the line}, whatever the innovation law.

\subsection{Factor graphs, and the tree}\label{sec:factor-graphs}

MRFs record \emph{which} variables interact but not \emph{how} the joint
factorises (a triangle of pairwise terms and one triple term have the same
graph). The representation that makes factorisation itself graphical ---
and on which message passing is most naturally defined --- is the
\emph{factor graph} \citep{kschischang2001factor,
koller2009probabilistic}.

\begin{definition}[Factor graph]
Given a factorisation $p(x) = \frac{1}{Z}\prod_{a} \psi_a(x_{\partial a})$,
its factor graph is the bipartite graph with a \emph{variable node} for
each $x_i$, a \emph{factor node} for each $\psi_a$, and an edge $(i, a)$
whenever $x_i \in x_{\partial a}$.
\end{definition}

The factor graph of \eqref{eq:bp-posterior} has $K$ variable nodes, $2K$
factor nodes ($1$ prior, $K{-}1$ transitions, $K$ observations --- the
data $x_k$ are constants clamped inside their factors, not nodes), and
$3K - 1$ edges: each transition factor has degree $2$, prior and
observations degree $1$. A connected graph with (nodes${}-1$) edges is a
\emph{tree}. Pictorially it is a \emph{caterpillar}: a spine of variables
and transition factors, one evidence leaf per vertebra --- in the words of
the supervision discussion, ``very tree-like: a straight line with things
on the side''. Every interior variable has degree three; every message
update combines a small, fixed number of inputs. Recording the conclusion
the rest of the chapter exploits:

\begin{proposition}\label{prop:chain-tree}
For any Markov-chain prior and any per-frame observation model, the
posterior $p(a \,|\, x)$ factorises over a tree-shaped factor graph, and
all $K$ smoothing marginals are computed exactly by $O(K)$ message
\emph{updates} (one forward and one backward sweep).
\end{proposition}

Note carefully what Proposition~\ref{prop:chain-tree} does \emph{not} say:
it does not say the total \emph{cost} is $O(K)$, because it does not say
the messages are finitely parametrisable. For discrete chains each message
is a probability vector and each update is finite arithmetic; for
continuous chains each message is a \emph{function}, each update is an
integral in function space, and the cost per update is unbounded in
general. The count of updates is linear in $K$; the price of an update
depends entirely on how messages are represented. Keeping
``tree-exactness of the algorithm'' separate from ``representability of
the messages'' is the conceptual hinge of this thesis: the first is a
graph property, settled here once and for all; the second is where all the
difficulty --- and all of Sections~\ref{sec:bp-gaussian} and
\ref{sec:bp-nongaussian} --- lives.

\section{The sum--product algorithm}\label{sec:bp-chain}

On a factor graph, sum--product belief propagation passes messages along
every edge \citep{pearl1988probabilistic, kschischang2001factor,
mezard2009information}: variable-to-factor messages multiply the incoming
factor messages from all \emph{other} neighbours, and factor-to-variable
messages integrate the factor against the incoming variable messages,
\begin{equation}
  \nu_{i \to a}(a_i) \propto
    \prod_{b \in \partial i \setminus a} \hat\nu_{b \to i}(a_i),
  \qquad
  \hat\nu_{a \to i}(a_i) \propto
    \int \psi_a(a_{\partial a})
    \prod_{j \in \partial a \setminus i} \nu_{j \to a}(a_j)\;
    \dd a_{\partial a \setminus i},
  \label{eq:bp-rules}
\end{equation}
with the belief at node $i$ proportional to
$\prod_{a \in \partial i} \hat\nu_{a \to i}(a_i)$. The ``exclude the
recipient'' rule is the heart of the algorithm: the message $i \to a$
summarises everything node $i$ knows \emph{except} what it learned from
$a$, so that no piece of evidence is ever counted twice along a path. On a
tree this bookkeeping is perfect: each edge separates the graph into two
disjoint halves, the message crossing it summarises exactly the half
behind it, and the equations \eqref{eq:bp-rules} have a unique solution,
reached in one inward and one outward pass, whose beliefs are the
\emph{exact} marginals \citep[Theorem~14.1]{mezard2009information}. On
graphs with cycles the same updates, iterated, define \emph{loopy} BP ---
extraordinarily useful, but approximate \citep{weiss2001correctness,
malioutov2006walk}; none of that is needed here, because our graph is a
tree.

For \emph{discrete} variables each message is a finite probability vector
and \eqref{eq:bp-rules} is arithmetic. For \emph{continuous} variables
each message is a probability \emph{density} --- an infinite-dimensional
object --- and the update integrals must be carried out in function space.
This is the obstruction, identified sharply in supervision, that makes
``BP on continuous chains'' a research question rather than a textbook
exercise. There are exactly three ways out: (i) discrete data (retreat to
finite alphabets); (ii) a family of densities closed under the updates ---
the Gaussian route, exact for our Gaussian chain, as we prove next;
(iii) numerical function representation --- the grid route, which
Section~\ref{sec:bp-nongaussian} uses as exact reference for non-Gaussian
chains.

\subsection{The chain sweeps, with the bookkeeping made explicit}

On the chain, \eqref{eq:bp-rules} collapses into one forward and one
backward sweep. The bookkeeping deserves unusual care --- an earlier
version of this project's notes fell into a trap here, caught only by the
numerical audit.

\begin{proposition}[Convention A]\label{prop:bp-convA}
Define messages
\begin{align}
  \mu_{\to 0}(a_0) &:= \psi_{\mathrm{pr}}(a_0),
  &
  \mu_{\to k+1}(a_{k+1}) &:= \int
    \psi^{(k)}_{\mathrm{tr}}(a_k, a_{k+1})\,
    \psi^{(k)}_{\mathrm{ob}}(a_k)\,
    \mu_{\to k}(a_k)\; \dd a_k ,
  \label{eq:bp-fwd}\\
  \mu_{\leftarrow K-1}(a_{K-1}) &:= 1,
  &
  \mu_{\leftarrow k-1}(a_{k-1}) &:= \int
    \psi^{(k-1)}_{\mathrm{tr}}(a_{k-1}, a_k)\,
    \psi^{(k)}_{\mathrm{ob}}(a_k)\,
    \mu_{\leftarrow k}(a_k)\; \dd a_k .
  \label{eq:bp-bwd}
\end{align}
Then for every $k$,
\begin{equation}
  p(a_k \,|\, x) \;\propto\;
  \mu_{\to k}(a_k)\;
  \psi^{(k)}_{\mathrm{ob}}(a_k)\;
  \mu_{\leftarrow k}(a_k) :
  \label{eq:bp-combine}
\end{equation}
the forward message into $a_k$ carries the evidence $x_0, \dots, x_{k-1}$
only, the backward message carries $x_{k+1}, \dots, x_{K-1}$ only, and the
local evidence enters once, at combination time.
\end{proposition}

The proof is direct marginalisation: integrating \eqref{eq:bp-posterior}
over all $a_j$, $j \neq k$, splits at $a_k$ by the Markov property into
independent left and right halves, and peeling the innermost variable of
each half yields the recursions. What deserves the reader's full attention
is the \emph{bookkeeping}:

\begin{remark}[The double-counting trap]\label{rmk:bp-trap}
A natural-looking alternative (call it Convention B) absorbs the local
evidence into the forward message,
$\mu^{\mathrm{B}}_{\to k} = \mu^{\mathrm{A}}_{\to k} \cdot
\psi^{(k)}_{\mathrm{ob}}$. Under Convention B the product
\eqref{eq:bp-combine} counts $x_k$ \emph{twice} and must be corrected by
dividing one factor back out. Convention A is the bookkeeping under which
the textbook combination rule is correct \emph{as written}. The audit was
decisive: under Convention A the BP score matches the matrix score to
$10^{-14}$; under naive Convention B the disagreement is order one. This
is the single most error-prone point of the construction --- hence every
message in this thesis is stated with its evidence content made explicit.
\end{remark}

Each sweep computes $K - 1$ messages, and the combination at each node
multiplies three univariate functions: $O(K)$ updates, as
Proposition~\ref{prop:chain-tree} promised --- with the cost of each
update still an open bill, to be settled family by family.

\section{The Gaussian case: closure, and the Kalman filter}
\label{sec:bp-gaussian}

For the Gaussian chain, two elementary facts make the message space
finite-dimensional --- and they are the \emph{only} facts needed.

\begin{toolbox}{The two closure lemmas of Gaussian message passing}
\emph{Lemma 1 (products).} For Gaussians in the same variable,
\begin{equation*}
  \Nd(a; m_1, v_1) \times \Nd(a; m_2, v_2)
  \;\propto\; \Nd(a; m, v),
  \qquad
  \frac{1}{v} = \frac{1}{v_1} + \frac{1}{v_2},
  \qquad
  \frac{m}{v} = \frac{m_1}{v_1} + \frac{m_2}{v_2} :
\end{equation*}
add the exponents; the sum of two quadratics in $a$ is a quadratic in $a$,
with precisions and precision-weighted means adding. (In the natural
parameters $\lambda = 1/v$, $\theta = m/v$ the rule is literally
$(\lambda, \theta) \mapsto (\lambda_1 + \lambda_2,\,
\theta_1 + \theta_2)$.)

\emph{Lemma 2 (affine marginalisation).} Pushing a Gaussian through the
AR(1) transition,
\begin{equation*}
  \int \Nd(a'; \alpha a, \sigma_\eta^2)\; \Nd(a; m, v)\; \dd a
  \;=\; \Nd(a';\, \alpha m,\, \alpha^2 v + \sigma_\eta^2),
\end{equation*}
because $a' = \alpha a + \eta$ with independent Gaussian summands: means
and variances transform as stated, and Gaussianity is preserved. The same
computation read in the reverse direction gives the backward analogue
$\Nd\big(a;\, m'/\alpha,\, (v' + \sigma_\eta^2)/\alpha^2\big)$.

Together: the Gaussian family is closed under exactly the two operations
the recursions \eqref{eq:bp-fwd}--\eqref{eq:bp-bwd} perform. Note also
that the local evidence, read as a function of $a_k$, is itself Gaussian
--- $\psi^{(k)}_{\mathrm{ob}}(a_k) \propto
\Nd(a_k;\, e^{t} x_k,\, \Deltat e^{2t})$, a pseudo-observation at the
deconvolved location $e^t x_k$ with precision $e^{-2t}/\Deltat$, the
pinning strength of Chapter~\ref{ch:gaussian}. \hfill$\square$
\end{toolbox}

\begin{theorem}[Gaussian closure of BP on the chain]\label{thm:bp-closure}
For the Gaussian AR(1) chain under the OU channel, every BP message and
every belief is exactly Gaussian, at every node, for every
$(K, \alpha, t)$. Each message is carried by two numbers, so each update
is $O(1)$ arithmetic and the two sweeps compute the exact posterior
marginals --- hence, via the score--posterior identity, the exact joint
score --- at total cost $O(K)$. \emph{Here, and only here, the linear
update count of Proposition~\ref{prop:chain-tree} becomes a linear cost.}
\end{theorem}

\begin{proof}[Derivation]
Induction along each sweep. The initial messages are Gaussian (the prior;
the flat message is the $\lambda \to 0$ limit of the family). Each step of
\eqref{eq:bp-fwd}--\eqref{eq:bp-bwd} is one application of Lemma~1 (absorb
the local evidence) followed by one of Lemma~2 (push through the
transition); the combination \eqref{eq:bp-combine} is Lemma~1 twice.
Exactness follows from tree-exactness of sum--product.
\end{proof}

\begin{remark}[Closure is algebra, not a central limit theorem]
\label{rmk:bp-no-clt}
The closure holds at node degree two just as it would at degree two
thousand: it is a property of the Gaussian \emph{family} under
multiplication and affine marginalisation, not an
aggregation-of-many-influences effect. This answers precisely the doubt
raised in supervision --- whether Gaussian messages on a chain could be
trusted, given that the usual central-limit justification needs many
incoming messages. For \emph{exact} Gaussian BP, no CLT is invoked
anywhere, so there is nothing to fail. The CLT enters only to justify the
\emph{per-node} closures of statistical physics (AMP/TAP) --- and there,
on the chain, it fails exactly as predicted; Appendix~\ref{app:amp} makes
that failure quantitative, complete with an exact breakdown time and the
critical coupling $\alpha_c = \sqrt{2} - 1$.
\end{remark}

\subsection{The sweeps are the Kalman filter and smoother}
\label{sec:bp-kalman}

Executing the Gaussian sweeps in mean--variance arithmetic reproduces,
line by line, the most famous algorithm of estimation theory. The
identification deserves to be displayed rather than asserted.

\begin{toolbox}{The Kalman filter, derived from the BP sweep}
Write the forward message in Convention A as
$\mu_{\to k} = \Nd(a_k;\, m_{k|k-1},\, P_{k|k-1})$ --- in filtering
language, the \emph{predicted} state given past observations. One BP step
does two things.

\emph{Update (absorb the evidence; Lemma 1).} Multiplying by the
pseudo-observation $\Nd(a_k; e^t x_k, \Deltat e^{2t})$ gives the
\emph{filtered} Gaussian $\Nd(a_k; m_{k|k}, P_{k|k})$ with
\begin{align*}
  \frac{1}{P_{k|k}} &= \frac{1}{P_{k|k-1}} + \frac{e^{-2t}}{\Deltat},\\
  m_{k|k} &= P_{k|k} \left( \frac{m_{k|k-1}}{P_{k|k-1}}
        + \frac{e^{-t} x_k}{\Deltat} \right)
  = m_{k|k-1} + G_k \big( e^{t} x_k - m_{k|k-1} \big),
\end{align*}
where the second form defines the \emph{Kalman gain}
$G_k = P_{k|k} \, e^{-2t}/\Deltat \in (0,1)$: the filtered mean is the
prediction, corrected by a gain-weighted innovation --- the celebrated
predictor--corrector structure \citep{kalman1960new}.

\emph{Predict (push through the dynamics; Lemma 2).} Propagating through
the AR(1) transition,
\begin{equation*}
  m_{k+1|k} = \alpha\, m_{k|k},
  \qquad
  P_{k+1|k} = \alpha^2 P_{k|k} + \sigma_\eta^2 :
\end{equation*}
contract the mean, inflate the variance by the innovation noise. The
forward BP sweep \emph{is} the Kalman filter --- not analogous to it,
equal to it, message by message.

\emph{Smoothing (combine both directions).} The belief
\eqref{eq:bp-combine} multiplies the filtered forward Gaussian by the
backward message (Lemma 1 once more), producing the posterior marginal
given \emph{all} observations, $\Nd(a_k; m_{k|K-1}, P_{k|K-1})$. Carrying
out the same algebra backwards recursion-by-recursion yields the classical
fixed-interval smoother of \citet{rauch1965maximum}: the two-sweep BP
schedule on the chain is exactly the Kalman filter followed by the RTS
smoother, the Gaussian specialisation of sum--product
\citep{kschischang2001factor, bishop2006pattern}. \hfill$\square$
\end{toolbox}

The score--posterior identity then gives the punchline of RQ2: \emph{the
$k$-th component of the joint diffusion score is an affine function of the
Kalman-smoothed estimate of frame $k$},
\begin{equation}
  S_k(x, t)
  \;=\; \frac{e^{-t}\, m_{k|K-1} - x_k}{\Deltat} .
  \label{eq:bp-score-kalman}
\end{equation}
Score computation for Markov sequences is a smoothing problem, and seventy
years of smoothing theory apply verbatim.

\paragraph{Machine-precision agreement.} The computational companion runs
BP against the matrix route of Chapter~\ref{ch:gaussian} over $80$
configurations --- $K \in \{2,3,5,8,16\}$, $\alpha \in \{0.2, 0.5, 0.9,
-0.4\}$, $t \in \{0.05, 0.3, 1, 3\}$, random $x$: maximum disagreement
$1.3 \times 10^{-15}$ on posterior means, $3.9 \times 10^{-15}$ on
variances, $1.2 \times 10^{-14}$ on scores --- double-precision rounding
accumulated along the recursion, nothing more. An independent
implementation, written from scratch against the supervision-call
specification alone, reproduces the same agreement; and a worked $K = 3$
instance with every message spelled out numerically is part of the
companion notebooks. One qualitative detail of that instance is worth
repeating: at $x = (1.2, -0.4, 0.8)$ the score component $S_1$ is
\emph{positive} although $x_1 < 0$ --- the neighbours pull the posterior
mean of the middle frame positive. No per-frame marginal score can produce
this; it is inter-frame evidence sharing, the joint effect of
Section~\ref{sec:g-K2}, in numbers.

\section{The Gaussian closure beyond Gaussianity}\label{sec:bp-nongaussian}

On the Laplace chain of Chapter~\ref{ch:laplace}, BP is still exact in
principle but its messages are genuinely infinite-dimensional
(Remark~\ref{rmk:l-breakdown}). This is where the closure question becomes
a measurement. Two solvers are compared on identical data.

\paragraph{Exact reference: grid BP.} Messages are represented by their
values on $M$ equispaced points in $[-A, A]$ and updated by trapezoidal
quadrature --- option (iii) of Section~\ref{sec:bp-chain}, at cost
$O(K M^2)$: the honest price of function-valued messages. Calibration on
the exactly solvable Gaussian chain shows the trapezoid rule on analytic,
rapidly decaying integrands converges \emph{spectrally} (aliasing error
$\sim e^{-2\pi^2 \sigma^2/\dd x^2}$), so failure is abrupt, not gradual:
at $M = 101$, $A = 8$ the score error sits at the $10^{-15}$ floor down to
$t = 0.02$ and fails only when the likelihood width $\sqrt{2t}$ approaches
the grid step. The working configuration for all non-Gaussian experiments
is $M = 401$, $A = 8$ (reference $M = 801$), with the resolution
constraint $\sqrt{2t} \gtrsim 3\,\dd x$ checked automatically; grid
self-convergence error remains $10^4$--$10^7$ times smaller than the
effect being measured at every $t$. Within its validated regime, grid BP
\emph{is} the exact score.

\paragraph{The object under test: Gaussian-projected BP.} The same
sweeps, with every message forced into the Gaussian family by moment
matching --- the algorithm one would actually deploy when quadrature is
unaffordable, and the chain analogue of the per-node closures of
Appendix~\ref{app:amp}.

\paragraph{The error instrument.} All score errors are amplified
posterior-mean errors: for any two solvers producing means
$\widehat m, m_{\mathrm{ref}}$ and scores by the score--posterior
identity,
\begin{equation}
  \widehat s - s_{\mathrm{ref}}
  = \frac{e^{-t}}{\Deltat}\,
    \big( \widehat m - m_{\mathrm{ref}} \big),
  \qquad
  \frac{e^{-t}}{\Deltat} \sim \frac{1}{2t} \;\; (t \downarrow 0),
  \label{eq:bp-master-error}
\end{equation}
an exact identity (algebra, not approximation) whose numerical residual is
logged in every experiment and sits at $10^{-15}$--$10^{-12}$ throughout.
Equation \eqref{eq:bp-master-error} explains \emph{a priori} why small
diffusion times are the danger zone: whatever the closure does to the
posterior mean is amplified by $1/2t$ in the score.

\subsection{Findings on the Laplace chain}

On the Laplace chain ($\rho = 0.85$, $n = 40$, $60$ trials per $t$,
reference $M = 801$): the median relative score error of the Gaussian
closure is
\begin{center}
\begin{tabular}{@{}lcccc@{}}
\toprule
$t$ & $0.02$ & $0.08$ & $0.4$ & $2.4$ \\
\midrule
median relative score error & $0.39$ & $0.17$ & $3.5 \times 10^{-2}$
& $7.6 \times 10^{-5}$ \\
\bottomrule
\end{tabular}
\end{center}
with failure probability $\Prob(\mathrm{err} > 0.1)$ equal to $1$ for
$t \leq 0.08$ and $0$ for $t \geq 0.9$
(Figure~\ref{fig:bp-closure}). The closure error is thus of order tens of
percent in the small-$t$ regime and decays below $10^{-3}$ once diffusion
noise Gaussianises the tilted posteriors --- the quantitative form of the
supervision intuition that ``at long diffusion times the world is
Gaussian, so probably you have no problem''. Inspection of the worst
trials shows the mechanism: the reference belief is sharply non-Gaussian
(occasionally bimodal), the single Gaussian commits to a compromise mean,
and \eqref{eq:bp-master-error} amplifies the shift. The score
\emph{direction}, however, remains far more accurate than its magnitude
throughout: the median cosine similarity between projected and reference
scores is $0.92$ even at $t = 0.02$ (where the magnitude error is $39\%$),
exceeds $0.99$ for all $t \geq 0.12$, and the worst single trial of the
whole sweep stays above $0.78$.

\begin{figure}[t]\centering
\includegraphics[width=0.74\linewidth]{figures/fig_laplace_closure.pdf}
\caption{\textbf{Gaussian closure error on the Laplace chain} (median and
10--90\% band) against the grid error budget measured on identical data
($\rho = 0.85$): the closure error dominates the reference error by
$10^4$--$10^7$ at every $t$, so the measurement is trustworthy where it is
used.}
\label{fig:bp-closure}
\end{figure}

\subsection{Varying the innovation: what makes closure hard}

Sweeping six innovation families with \emph{identical} covariance
$\rho^{|i-j|}$ --- Laplace, Student-$t$ ($\nu = 5, 8$), and symmetric
Gaussian mixtures of increasing separation --- isolates the shape of the
innovation law as the only moving part. Three findings:
\begin{enumerate}
  \item \emph{Noise level dominates.} Every family's error decays to
  $\sim 10^{-3}$ or below by $t = 2.4$; all meaningful closure error lives
  at $t \lesssim 0.4$.
  \item \emph{Both signs of non-Gaussianity hurt; bimodality hurts most.}
  At $t = 0.05$, $\rho = 0.5$: strongly bimodal mixture $0.94$, Laplace
  $0.52$, Student-$t_{\nu=5}$ $0.32$, mildly bimodal mixture $0.07$. A
  single Gaussian cannot represent two posterior modes, and commits to the
  gap between them.
  \item \emph{Stronger correlation helps, not hurts.} Laplace at
  $t = 0.05$: error $0.43, 0.27, 0.13$ for $\rho = 0.5, 0.85, 0.95$.
  Informative neighbours act like extra Gaussian observations and
  Gaussianise the local tilted posterior; the naive intuition that message
  errors ``propagate farther'' at strong coupling is not what happens on
  chains.
\end{enumerate}

\subsection{Truncating the matrix versus truncating the inference}
\label{sec:bp-truncation}

The last supervision question --- take the exact score and truncate it to
tridiagonal: how bad? --- admits two inequivalent readings, and keeping
them apart turns out to matter. \emph{Truncate the answer} (variant a):
band the matrix, $\widehat S = -Q_t^{(b)} x$ with entries beyond
$|i - j| \leq b$ zeroed. \emph{Truncate the algorithm} (variant b): let
messages travel at most $r$ hops --- exact inference on the window
$x_{k-r}, \dots, x_{k+r}$, the local estimator of
Section~\ref{sec:g-locality}. Both estimators are linear in $x$, so their
RMS errors over the data law are deterministic traces --- no sampling.
At equal locality budget ($b = r = 1$, $K = 40$, $\alpha = 0.8$):
\begin{center}
\begin{tabular}{@{}lcccc@{}}
\toprule
relative RMS score error & $t = 0.05$ & $t = 0.3$ & $t = 1$ & $t = 3$ \\
\midrule
banded matrix ($b = 1$)      & $0.210$ & $0.302$ & $0.135$ & $0.0035$ \\
truncated inference ($r = 1$) & $0.123$ & $0.235$ & $0.130$ & $0.0035$ \\
\bottomrule
\end{tabular}
\end{center}
Two lessons. First, the cost of locality is non-monotone in $t$ and peaks
at intermediate times --- the same mid-diffusion window where the band of
$Q_t$ is fullest (Chapter~\ref{ch:gaussian}) and where generative
sampling spends its critical steps. Second, at small $t$ \emph{cutting the
algorithm beats cutting the answer} ($12.3\%$ vs $21.0\%$ at $t = 0.05$):
the windowed inference solves a correctly normalised problem on its
window, while zeroing entries of an inverse produces a matrix that is the
inverse of nothing and misstates the diagonal. The two coincide at large
$t$, where all off-diagonal content dies. The error of either truncation
decays geometrically in the allowed range, at the rate $q(\alpha, t)$ of
Theorem~\ref{thm:g-locality}.

\section{What this chapter established}\label{sec:bp-summary}

\emph{Exact} (proved, and verified to $10^{-12}$--$10^{-15}$): the factor
graph of the noisy chain is a tree, so the two BP sweeps compute all
smoothing marginals exactly in $O(K)$ message updates; on the Gaussian
chain the messages close algebraically in the Gaussian family --- no CLT
--- so the updates cost $O(1)$ each and BP, which is exactly the Kalman
filter plus RTS smoother, reproduces the matrix score to machine
precision. \emph{Numerical} (measured against validated exact
references): the Gaussian closure error on non-Gaussian chains ---
$\sim\!20$--$40\%$ at small $t$, $< 10^{-3}$ at large $t$, worst for
bimodal innovations, decreasing in $\rho$ --- with the score direction
much more robust than its magnitude; and the locality comparison, where
truncated inference beats the truncated matrix at small $t$.
\emph{Heuristic}: the architecture reading --- local, convolution-like
score networks with receptive field sized by $q(\alpha, t)$ (or grown
with $t$) are near-optimal outside the mid-diffusion window, and the
exactly quantified penalty inside it is the price of locality, not an
artefact. The per-node closures of statistical physics, and the exact
phase boundary they develop on this very model, are worked out in
Appendix~\ref{app:amp}; the findings, limitations, and follow-up
programme are consolidated in the next chapter.
